\documentclass{report}
\usepackage{amsfonts}
\newcommand\R{{\mathbb R}}
\newcommand\Q{{\mathbb Q}}
\newcommand\N{{\mathbb N}}
\pagestyle{empty}
\begin{document}
\large
\centerline{\Huge Analisi Matematica}
\renewcommand{\thefootnote}{ } 
\footnote{\sl Avete 3 ore di tempo. Ogni esercizio
vale sei punti. La sufficienza si ottiene con un punteggio $\geq$
18. Solo le risposte chiaramente giustificate saranno prese in
considerazione.}
\renewcommand{\thefootnote}{\arabic{footnote}} 
\vskip.1cm
\centerline{\Large Seconda Prova di Autovalutazione}
\vskip1cm
\begin{enumerate}
\item Dimostrare per induzione la seguente formula:\newline
per ogni $a>0$ ed
$n\in\N$
\[
(1+a)^n\geq 1+na+\frac {n(n-1)}2 a^2.
\]

\item Mettere in ordine crescente i seguenti numeri
\[
1000000^8;\quad 100!;\quad 10^{211};\quad 3^{4^9}.
\]

\item Calcolare
\[
\lim_{n\to\infty}e^nn^{-n}.
\]
\item Data la funzione $f:\R\to\R$ definita da $f(x):=3+e^{x-4}$
determinare la funzione inversa $f^{-1}$.
\item Calcolare $\ln 8$ con una cifra
decimale esatta. \`E permesso usare un calcolatore solo per estrarre
radici quadrate.

\item Discutere la convergenza di
\[
\sum_{i=1}^\infty \frac 1{n\ln n}.
\]



\end{enumerate}
\end{document}
